% ===================== Backup =====================
\section*{Backup}

\begin{frame}
  \centering \Huge \textbf{Backup slides}
\end{frame}

% B1
\begin{frame}{Backup: lattice ensemble and $\zetahat$ at $P_L=-1$}
  \begin{columns}[t]
    \begin{column}{0.48\textwidth}
      \begin{block}{From the thesis}
        \begin{itemize}
          \item $n_f=2+1$ clover-improved Wilson fermions (Yoon \emph{et al.} 2017)
          \item $\hat L^3\times\hat T = 32^3\times 96$, \ $a=0.11403$ fm, \ $L\approx3.65$ fm
          \item $\mN = 0.6228(28)/a = 1077.8(48)$ MeV
          \item $P_1=-2\pi/L=-339.8$ MeV
          \item Isovector $u-d$, connected contractions
        \end{itemize}
      \end{block}
      \begin{alertblock}{Not stated in the thesis}
        \todo{fill in: $m_\pi$, number of configurations and Jackknife bins, $t_{\rm snk}-t_{\rm src}$, smearing}
      \end{alertblock}
    \end{column}
    \begin{column}{0.50\textwidth}
      Setting the bracket of the constraint to zero:
      \[ \zetahat =\sqrt{\frac{P_{1}^2}{\mN^{2}\left(\frac{\mN^{2}}{P_{1}^{2}}+2\right)}} = \pm 0.090772 \]
      \[ \zetahat=\frac{-v_{1}P_{1}}{\sqrt{v_{1}^{2}+v_{3}^{2}}\ \mN} \ \Rightarrow\ v_{1}=\pm 0.301\,v_{3} \]
    \end{column}
  \end{columns}
\end{frame}

% B2
\begin{frame}{Backup: coverage in $(b_L, b_T)$ and $(P\cdot b,\ b^2)$}
  \begin{columns}[c]
    \begin{column}{0.48\textwidth}
      \centering
      \includegraphics[width=0.8\linewidth]{bL_bT_plot.pdf}\\
      {\scriptsize Lattice $\vect b$ vectors in the $(b_L,b_T)$ plane}
    \end{column}
    \begin{column}{0.48\textwidth}
      \centering
      \includegraphics[width=\linewidth]{Pb_bsq_plot.pdf}\\
      {\scriptsize Coverage in $(P\cdot b,\ b^2)$ for the four momenta}
    \end{column}
  \end{columns}
  \begin{itemize}
    \item $|P\cdot b|\le|\vect P|\,|\vect b|$: the reach in $P\cdot b$ grows with $|P_L|$.
    \item Small $x$ is controlled by large $P\cdot b$, outside the data: the least constrained region.
  \end{itemize}
\end{frame}

% B3
\begin{frame}{Backup: the $|\vect b|\ge3a$ cut for $\tAmp_{12B}^{\text{Re}}$}
  \vspace{2mm}
  \centering
  \includegraphics[width=\linewidth]{A12B_Re_b_lessthan3_and_ge3.pdf}
\end{frame}

% B3b
\begin{frame}{Backup: $\tAmp_{12B}^{\text{Im}}$ is consistent with zero}
  \centering
  \includegraphics[width=\linewidth]{ImA12B_Pb_bsq_eta8_PL-3_2D.pdf}\\
  {\scriptsize $\tAmp_{12B}^{\text{Im}}$ at $P_L=-3$, $\eta|v|/a=8$ (all $b^2$, including $|\vect b|<3a$): consistent with zero, so dropped.}
\end{frame}

% B4
\begin{frame}{Backup: PySR benchmark fits}
  \begin{columns}[c]
    \begin{column}{0.62\textwidth}
      \centering
      \includegraphics[height=0.8\textheight, width=\linewidth, keepaspectratio]{01_A2B_PySR_Benchmark_Fits.pdf}
    \end{column}
    \begin{column}{0.36\textwidth}
      \begin{itemize}
        \item Run on $\tAmp_{2B}^{\text{Re}}$ at $P_L=-1$ only.
        \item $\mathcal{P}_{\text{asymptotic}}$ sums $|\tAmp^{\text{model}}|$ at large $b_L$.
        \item Every formula with spatial dependence depends on $b_L, b_T$ only through $b^2$: a guide, not the final model.
      \end{itemize}
    \end{column}
  \end{columns}
\end{frame}

% B5
\begin{frame}{Backup: $c(\eta)$ candidates and the Fourier-transform width}
  \centering
  \includegraphics[height=0.82\textheight, keepaspectratio]{A2BfitFT_eta_Asymptotic_P1.pdf}
\end{frame}

% B5b
\begin{frame}{Backup: $\eta$-independence of $d$, and the $c(\eta)$ choice}
  \begin{columns}[c]
    \begin{column}{0.48\textwidth}
      \centering
      \includegraphics[width=0.9\linewidth]{A2B_d_eta_Independence_P1.pdf}\\
      {\scriptsize $d(\eta)$ for $\tAmp_{2B}^{\text{Re}}$: slope $0.00122(87)$}
    \end{column}
    \begin{column}{0.50\textwidth}
      \begin{itemize}
        \item The $x$-width of $\tAmp_{2B}^{\text{Re}}$ moves with the choice of $c(\eta)$ (previous slide).
        \item $c_2(\eta)=c\big(1-\tfrac{k_1}{\eta}-\tfrac{k_2}{\eta^2}\big)$: lowest AIC, suppresses antiquarks at $x\lesssim-0.2$, consistent with the $x$-integrated results.
        \item The $b_L^2$ coefficient is $d+c_2(\eta)$, which stays positive over $\eta\in[6,10]$.
        \item \todo{how $c(\eta)$ with $k_1,k_2$ enters the transform and the $\zetahat$ extrapolation}
      \end{itemize}
    \end{column}
  \end{columns}
\end{frame}

% B6
\begin{frame}{Backup: integrability of the power-law models}
  \begin{itemize}
    \item Even amplitudes: $\big[1+C\,(P\cdot b)^2+\ldots\big]^{-j}\sim|P\cdot b|^{-2j}$ $\Rightarrow$ absolutely integrable for $j>1/2$.
    \item $\tAmp_{2B}^{\text{Im}}\propto b_L\,[\ldots]^{-j}\sim|P\cdot b|^{1-2j}$ $\Rightarrow$ absolutely integrable only for $j>1$.
    \item $j_{\text{im}}$ at the simulated momenta $P_L=-1,-2,-3,-4$: \ $1.108(21)$, $1.141(20)$, $1.057(15)$, $1.0084(79)$, all $>1$.
    \item Extrapolated to $\zetahat\to\infty$:
  \end{itemize}
  \[ j_{\text{re}} = 2.0667(69), \quad j_{\text{im}} = 0.9921(92), \quad j_{\text{reA12B}} = 2.0026(40) \]
  \begin{itemize}
    \item At $j_{\text{im}}\approx1$ the sine transform converges only conditionally; this is the sharp structure of $\tAmp_{2B}^{[x]\text{Im}}$ near $x=0$.
    \item Power laws are not entire functions, so support extends beyond $|x|\le1$; numerically $\int_{-\infty}^{\infty}$ and $\int_{-1}^{1}$ agree within errors (next slide).
  \end{itemize}
\end{frame}

% B7
\begin{frame}{Backup: $x$-integrals at $\zetahat\to\infty$, $|\vprp b|=3a$}
  \centering
  \renewcommand{\arraystretch}{1.5}
  {\small
  \begin{tabular}{lcccc}
    \toprule
    Integral & $[-\infty, \infty]$ & $[-1, 1]$ & $[-1, 0]$ & $[0, 1]$ \\
    \midrule
    $\int dx\, \tAmp^{[x]\text{Re}}_{2B}$ & $0.751(79)$ & $0.751(80)$ & $0.375(40)$ & $0.375(40)$ \\
    $\int dx\, \tAmp^{[x]\text{Im}}_{2B}$ & $6(96)\times10^{-17}$ & $3(24)\times10^{-16}$ & $0.165(27)$ & $-0.169(29)$ \\
    $\int dx\, \tilde{f}_1^{(0)}$ & $1.50(16)$ & $1.50(16)$ & $0.421(81)$ & $1.09(11)$  \\
    $\int dx\, \tilde{f}_{1T}^{\perp(1)}$ & $0.302(53)$ & $0.302(53)$ & $0.151(27)$ & $0.151(27)$ \\
    $\mN \int dx\, \tilde{f}_{1T}^{\perp(1)} / \int dx\, \tilde{f}_1^{(0)}$ & $0.217(38)$ & $0.217(38)$ & $0.387(89)$ & $0.150(26)$ \\
    \bottomrule
  \end{tabular}}
  \raggedright
  \begin{itemize}
    \item First four rows are unrenormalized; only the ratio (last row) is quoted as a physical quantity.
    \item $\int_{-\infty}^{\infty}dx$ of the transform equals the amplitude at $\bvec\tcdot P=0$ (an identity). The content is the split into $[-1,0]$ and $[0,1]$.
  \end{itemize}
\end{frame}

% B8
\begin{frame}{Backup: all twelve parameter extrapolations}
  \centering
  \includegraphics[height=0.8\textheight, keepaspectratio]{Param_Extrapolations_1byZetahat_sq_P234_bmin3_eta610_b29.0.pdf}\\
  {\scriptsize Rows: $\tAmp_{2B}^{\text{Re}}$, $\tAmp_{2B}^{\text{Im}}$, $\tAmp_{12B}^{\text{Re}}$. Each fit: 3 points, 2 parameters.}
\end{frame}

% B9
\begin{frame}{Backup: Sivers shift at the four momenta vs $1/\zetahat$}
  \centering
  \includegraphics[width=0.9\linewidth]{SimulFit_Distinct_1_over_PL_bmin3_eta610_b29.0_Sivers.pdf}
\end{frame}

% B10
\begin{frame}{Backup: extrapolated amplitudes and TMDs}
  \centering
  \includegraphics[height=0.82\textheight, keepaspectratio]{Extrapolated_AmpLevel_P234_Sivers_1_over_zetahat_sq_bmin3_eta610_b29.0_Amps.pdf}
\end{frame}

% B11
\begin{frame}{Backup: full $x$-dependence at $P_L=-3$}
  \centering
  \includegraphics[width=0.75\linewidth]{FT_xdep_A2B_A12B_bmin3_eta610_b2-9_PL-3.pdf}\\
  {\scriptsize Top: $\tAmp^{[x]\text{Re}}_{2B}$, $\tAmp^{[x]\text{Im}}_{2B}$, $\tilde f_1^{(0)}$. Bottom: $\tilde f_{1T}^{\perp(1)}$, $\langle k_y\rangle_{TU}$, $x\langle k_y\rangle_{TU}$.}
\end{frame}

% B12
\begin{frame}{Backup: global fit at $P_L=-3$}
  \centering
  \includegraphics[height=0.82\textheight, keepaspectratio]{Combined_Pb_bsq_eta8_PL-3_WithFit_2D.pdf}
\end{frame}

% B12b
\begin{frame}{Backup: separable models vs the coupled power law}
  \centering
  {\scriptsize
  \renewcommand{\arraystretch}{1.3}\setlength{\tabcolsep}{3.5pt}
  \begin{tabular}{lccccccc}
    \toprule
    Model: longitudinal $\times$ transverse & \multicolumn{4}{c}{$\chi^2/\text{dof}$ at $P_L=-1,-2,-3,-4$} & $c^{\text{Im}}_{2B}$ & $c^{\text{Re}}_{12B}$ & $a^{\text{Im}}_{2B}$ \\
    \midrule
    $e^{-c|P\cdot b|}\times e^{-d|\vect b|}$ & $1.406(92)$ & $0.956(75)$ & $0.622(61)$ & $0.455(53)$ & $0.0010(262)$ & $0.045(23)$ & $-0.167(16)$ \\
    $e^{-c|P\cdot b|}\times(1+d\,b^2)^{-j}$ & $1.286(89)$ & $0.938(75)$ & $0.623(61)$ & $0.447(53)$ & $0.0045(248)$ & $0.028(24)$ & $-0.167(31)$ \\
    $e^{-c(P\cdot b)^2}\times e^{-d|\vect b|}$ & $2.389(121)$ & $2.005(111)$ & $0.945(76)$ & $0.453(52)$ & $0.568(171)$ & $0.0038(69)$ & $+0.050(121)$ \\
    $e^{-c(P\cdot b)^2}\times(1+d\,b^2)^{-j}$ & $2.316(120)$ & $1.981(111)$ & $0.948(76)$ & $0.448(52)$ & $0.602(48)$ & $0.0016(61)$ & $+0.168(92)$ \\
    \midrule
    \textbf{coupled power law in $(P\cdot b, b^2)$} & $1.85(10)$ & $1.197(85)$ & $0.712(66)$ & $0.434(51)$ & $\mathbf{0.2253(69)}$ & $\mathbf{0.0773(46)}$ & $\mathbf{+0.267(13)}$ \\
    \bottomrule
  \end{tabular}}
  \raggedright
  \begin{itemize}
    \item Parameters shown at $P_L=-1$. Separable models: $c^{\text{Im}}_{2B}$ or $c^{\text{Re}}_{12B}$ consistent with zero, $a^{\text{Im}}_{2B}$ sign flips, or poor $\chi^2$.
  \end{itemize}
\end{frame}

% B13
\begin{frame}{Backup: factorization ratio at $\zetahat\to\infty$}
  \centering
  \includegraphics[height=0.68\textheight, keepaspectratio]{Factorization_Ratio_vs_x_4curves_P234_bmin3_eta610.pdf}\\
  {\scriptsize $R_b(x)=\langle k_y\rangle_{TU}(x,|\vprp b|)/\langle k_y\rangle_{TU}(x,0.34\,\text{fm})$: not flat in $x$}
  \begin{itemize}
    \item This tests the extrapolated fitted model; the raw-data ratio test (next slide) is too sparse to decide on its own.
  \end{itemize}
\end{frame}

% B13a
\begin{frame}{Backup: ratio test directly on lattice data ($P_L=-1$)}
  \centering
  \includegraphics[height=0.82\textheight, keepaspectratio]{Factorization_ratios_Pb034_lin_vs_quad_eta8_PL-1.pdf}
\end{frame}

% B13b
\begin{frame}{Backup: optimal separable factors and reconstruction}
  \centering
  \includegraphics[width=\linewidth]{Factorization_gh_and_Reconstruction_bmin3_eta610.pdf}\\
  {\scriptsize Left: $g(x)$; middle: $h(|\vprp b|)$ for four $x$-windows; right: data vs $g(x)h(|\vprp b|)$ at $|\vprp b|=0.34$ and $0.94$ fm}
\end{frame}

% B14
\begin{frame}{Backup: soft factor and renormalization}
  \begin{itemize}
    \item At fixed geometry: under $\Gamma$-independent multiplicative renormalization, the soft factor and $Z$ cancel in $\tAmp_{12B}/\tAmp_{2B}$ at the same $(\bvec, P, \eta v)$.
    \item This work: ratios of amplitudes at fixed geometry, with $\bvec\tcdot P$ scanned and Fourier transformed.
    \item The $(-\infty,\infty)$ entry of the last row of the integral table equals the $\bvec\tcdot P=0$ ratio (an identity).
    \item \todo{prepare the oral answer on what changes with $b_L$ at fixed $|\vprp b|$ (link length vs cusp angles)}
  \end{itemize}
\end{frame}

% B15
\begin{frame}{Backup: generalized shifts of other TMDs}
  \small
  \begin{align*}
    \langle \vect{k}_y \rangle_{UT}(\vprp{\bvec}^2,x)
      &= \mN \frac{\int d(\bvec\tcdot P)\, e^{ix(\bvec\tcdot P)}\, \tAmp_{4B}}{\int d(\bvec\tcdot P)\, e^{ix(\bvec\tcdot P)}\, \tAmp_{2B}} & &\text{(Boer--Mulders, T-odd)} \\
    \langle \vect{k}_x \rangle_{TL}(\vprp{\bvec}^2,x)
      &= -\mN \frac{\int d(\bvec\tcdot P)\, e^{ix(\bvec\tcdot P)}\, \tAmp_{7B}}{\int d(\bvec\tcdot P)\, e^{ix(\bvec\tcdot P)}\, \tAmp_{2B}} & &\text{(worm-gear $g_{1T}$, T-even)} \\
    \frac{\tilde h_{1}^{(0)}}{\tilde f_1^{(0)}}(\vprp{\bvec}^2,x)
      &= -\frac{\int d(\bvec\tcdot P)\, e^{ix(\bvec\tcdot P)}\, \tAmp_{9Bm}}{\int d(\bvec\tcdot P)\, e^{ix(\bvec\tcdot P)}\, \tAmp_{2B}} & &\text{(tensor charge, T-even)}
  \end{align*}
  \normalsize
  \begin{itemize}
    \item Unlike the $x$-integrated shifts ($\bvec\tcdot P=0$ only), each needs a scan in $\bvec\tcdot P$ under the same kinematic constraint.
  \end{itemize}
\end{frame}
